\subsection{Observed-data diagnostics and their target population}
\label{app:natural-diagnostics}
Theoretical selection uses conditional loss and exposure; naturally hidden demand does not supply either without assumptions. Our observed-sales-and-capacity negative-binomial experiment illustrates the issue. Among adjusted paired item-bootstrap comparisons, shape $\kappa=3$ is the only candidate whose validation likelihood difference from selected $\kappa=2$ includes zero, yet it breaks both realized calibration caps that $\kappa=2$ passes. The selected model overestimates aggregate error mass by 20.59\%; its pooled point pass does not validate conditional moments. Appendix~\ref{app:observable} provides the likelihood, dispersion grid, and count-matched supervised control.

Further forecasting experiments on naturally stockout-annotated FreshRetailNet test observable sales outcomes, not hidden demand. A forward-time residual learner reduces WAPE on the 8,020 non-stockout test rows from 0.32928 to 0.29201, and on all 14,000 recorded-sales rows from 0.36162 to 0.33401. These later comparisons reuse the previously opened seven-day cohort; most improvement remains without the risk feature (non-stockout WAPE 0.29289). They demonstrate useful forecast adaptation within that evaluated population, but cannot certify an acceptance constraint for demand coverage on unobserved stockout demand. Appendix~\ref{app:accuracy} records the controls and uncertainty. Thus accuracy improvement, acceptance utility, and identification are three distinct requirements of a forecasting engine.

